% GENERATED FILE. Edit canonical lessons, then run npm run generate:books.
% canonical-source-hash: fnv1a64:abc76f3078d89aa3
% canonical-lessons: MR-C482-dokedukhi, MR-C482-dukhne, MR-C482-kal-yene, MR-C482-kanhne, MR-C482-vivhalne

\chapter{Aches and Pains}
\label{ch:mr-aches-and-pains}
\hlchaptermodality{482}

\begin{chapteropening}
\textbf{By the end of this chapter:} \emph{I can say what aches and describe groaning in pain.}

Complete the last lesson of chapter 482: I can say what aches and describe groaning in pain.
\end{chapteropening}

\section[\texorpdfstring{\mr{डोकेदुखी} (ḍokedukhī)}{डोकेदुखी (ḍokedukhī)}]{\mr{डोकेदुखी} (ḍokedukhī) — a headache}

\label{lesson:MR-C482-dokedukhi}



\begin{quote}
Before the new one: say the Marathi for to blow one's nose, then the Marathi for to clear one's throat.
\end{quote}

\subsection*{You'll want to know: \mr{डोकेदुखी}}
\textbf{\mr{डोकेदुखी}} — \emph{ḍokedukhī} — \textquotedblleft{}a headache\textquotedblright{}.

\begin{quote}
\mr{उन्हात} \mr{फिरल्यामुळे} \mr{मला} \mr{डोकेदुखी} \mr{झाली}.
\end{quote}

\begin{quote}
\emph{unhāt phirlyāmuḷe malā ḍokedukhī jhālī.}
\end{quote}

\begin{quote}
I got a headache from walking about in the sun.
\end{quote}

Feminine.

\mr{डोके} + \mr{दुखी}, aching. Also, figuratively, a nuisance.

\begin{cousinweb}
\textbf{\mr{डोकेदुखी}} (\emph{ḍokedukhī}) is a Marathi compound of \mr{डोके}, \textquotedblleft{}head\textquotedblright{}, and \mr{दुखी}, \textquotedblleft{}aching\textquotedblright{}, from the verb \mr{दुखणे}, a tadbhava of Sanskrit \emph{duḥkhati} (CDIAL 6376).
\end{cousinweb}

\begin{grammarlens}[title={What you've built}]
One more word for describing what you lived and felt.
\end{grammarlens}

\subsection*{Your turn}
\begin{itemize}
\raggedright
  \item \emph{Say it:} \emph{ḍokedukhī}
  \item \emph{Say it:} \emph{ḍokedukhī}, once more
  \item \emph{Say it:} \emph{unhāt phirlyāmuḷe malā ḍokedukhī jhālī.} — I got a headache from walking about in the sun.
  \item \emph{Recall:} say the Marathi for to blow one's nose, then the Marathi for to clear one's throat, then say \emph{ḍokedukhī} again
\end{itemize}

\begin{tcolorbox}[breakable,colback=teal!4,colframe=teal!35!black,title={Before you move on}]
What does \mr{डोकेदुखी} mean? (\textquotedblleft{}A headache\textquotedblright{}.) Say it once more.
\end{tcolorbox}
\section[\texorpdfstring{\mr{दुखणे} (dukhṇe)}{दुखणे (dukhṇe)}]{\mr{दुखणे} (dukhṇe) — to ache; to hurt}

\label{lesson:MR-C482-dukhne}



\begin{quote}
Before the new one: say the Marathi for to clear one's throat, then the Marathi for a headache.
\end{quote}

\subsection*{You'll want to know: \mr{दुखणे}}
\textbf{\mr{दुखणे}} — \emph{dukhṇe} — \textquotedblleft{}to ache; to hurt\textquotedblright{}.

\begin{quote}
\mr{जास्त} \mr{वेळ} \mr{बसल्यामुळे} \mr{माझी} \mr{पाठ} \mr{दुखते}.
\end{quote}

\begin{quote}
\emph{jāsta veḷ baslyāmuḷe mājhī pāṭh dukhte.}
\end{quote}

\begin{quote}
My back aches from sitting too long.
\end{quote}

Past (to him …): \mr{त्याला} \mr{दुखलं}, \emph{tyālā dukhlaṁ}.

The aching part is the subject: \mr{माझं} \mr{डोकं} \mr{दुखतं}. With the dative and no part named: \mr{मला} \mr{खूप} \mr{दुखलं}, it hurt me a lot. \mr{दुखणं}, an ailment.

\begin{cousinweb}
\textbf{\mr{दुखणे}} (\emph{dukhṇe}) is a tadbhava, inherited from Sanskrit \emph{duḥkhati}, \textquotedblleft{}pains\textquotedblright{}; CDIAL lists Marathi \emph{dukhṇẽ} (CDIAL 6376).
\end{cousinweb}

\begin{grammarlens}[title={What you've built}]
One more word for describing what you lived and felt.
\end{grammarlens}

\subsection*{Your turn}
\begin{itemize}
\raggedright
  \item \emph{Say it:} \emph{dukhṇe}
  \item \emph{Say it:} \emph{dukhṇe}, once more
  \item \emph{Say it:} \emph{jāsta veḷ baslyāmuḷe mājhī pāṭh dukhte.} — My back aches from sitting too long.
  \item \emph{Recall:} say the Marathi for to clear one's throat, then the Marathi for a headache, then say \emph{dukhṇe} again
\end{itemize}

\begin{tcolorbox}[breakable,colback=teal!4,colframe=teal!35!black,title={Before you move on}]
What does \mr{दुखणे} mean? (\textquotedblleft{}To ache; to hurt\textquotedblright{}.) Say it once more.
\end{tcolorbox}
\section[\texorpdfstring{\mr{कळ} \mr{येणे} (kaḷ yeṇe)}{कळ येणे (kaḷ yeṇe)}]{\mr{कळ} \mr{येणे} (kaḷ yeṇe) — to get a twinge}

\label{lesson:MR-C482-kal-yene}



\begin{quote}
Before the new one: say the Marathi for a headache, then the Marathi for to ache.
\end{quote}

\subsection*{You'll want to know: \mr{कळ} \mr{येणे}}
\textbf{\mr{कळ} \mr{येणे}} — \emph{kaḷ yeṇe} — \textquotedblleft{}to get a twinge\textquotedblright{}.

\begin{quote}
\mr{जिना} \mr{चढताना} \mr{माझ्या} \mr{गुडघ्यात} \mr{कळ} \mr{आली}.
\end{quote}

\begin{quote}
\emph{jinā caḍhtānā mājhyā guḍghyāt kaḷ ālī.}
\end{quote}

\begin{quote}
I felt a twinge in my knee climbing the stairs.
\end{quote}

Past (to him …): \mr{त्याला} \mr{कळ} \mr{आली}, \emph{tyālā kaḷ ālī}.

\mr{कळ}, a sharp twinge: \mr{पोटात} \mr{कळ} \mr{येणे}, a cramp in the stomach. Dative: \mr{मला} \mr{कळ} \mr{आली}.

\begin{cousinweb}
\textbf{\mr{कळ} \mr{येणे}} (\emph{kaḷ yeṇe}) is a Marathi phrase built on \mr{कळ}, \textquotedblleft{}a sharp twinge\textquotedblright{}, whose origin is not settled, with \mr{येणे}, \textquotedblleft{}to come\textquotedblright{}.
\end{cousinweb}

\begin{grammarlens}[title={What you've built}]
One more word for describing what you lived and felt.
\end{grammarlens}

\subsection*{Your turn}
\begin{itemize}
\raggedright
  \item \emph{Say it:} \emph{kaḷ yeṇe}
  \item \emph{Say it:} \emph{kaḷ yeṇe}, once more
  \item \emph{Say it:} \emph{jinā caḍhtānā mājhyā guḍghyāt kaḷ ālī.} — I felt a twinge in my knee climbing the stairs.
  \item \emph{Recall:} say the Marathi for a headache, then the Marathi for to ache, then say \emph{kaḷ yeṇe} again
\end{itemize}

\begin{tcolorbox}[breakable,colback=teal!4,colframe=teal!35!black,title={Before you move on}]
What does \mr{कळ} \mr{येणे} mean? (\textquotedblleft{}To get a twinge\textquotedblright{}.) Say it once more.
\end{tcolorbox}
\section[\texorpdfstring{\mr{कण्हणे} (kaṇhṇe)}{कण्हणे (kaṇhṇe)}]{\mr{कण्हणे} (kaṇhṇe) — to groan}

\label{lesson:MR-C482-kanhne}



\begin{quote}
Before the new one: say the Marathi for to ache, then the Marathi for to get a twinge.
\end{quote}

\subsection*{You'll want to know: \mr{कण्हणे}}
\textbf{\mr{कण्हणे}} — \emph{kaṇhṇe} — \textquotedblleft{}to groan\textquotedblright{}.

\begin{quote}
\mr{तापामुळे} \mr{तो} \mr{रात्रभर} \mr{कण्हत} \mr{होता}.
\end{quote}

\begin{quote}
\emph{tāpāmuḷe to rātrabhar kaṇhat hotā.}
\end{quote}

\begin{quote}
He groaned all night with fever.
\end{quote}

Past (he …): \mr{तो} \mr{कण्हला}, \emph{to kaṇhlā}.

A low sound of pain.

\begin{cousinweb}
\textbf{\mr{कण्हणे}} (\emph{kaṇhṇe}) is a tadbhava: CDIAL lists Marathi \emph{kanhṇẽ}, \textquotedblleft{}to groan\textquotedblright{}, under Sanskrit \emph{kvaṇati}, \textquotedblleft{}sounds\textquotedblright{} (CDIAL 3632).
\end{cousinweb}

\begin{grammarlens}[title={What you've built}]
One more word for describing what you lived and felt.
\end{grammarlens}

\subsection*{Your turn}
\begin{itemize}
\raggedright
  \item \emph{Say it:} \emph{kaṇhṇe}
  \item \emph{Say it:} \emph{kaṇhṇe}, once more
  \item \emph{Say it:} \emph{tāpāmuḷe to rātrabhar kaṇhat hotā.} — He groaned all night with fever.
  \item \emph{Recall:} say the Marathi for to ache, then the Marathi for to get a twinge, then say \emph{kaṇhṇe} again
\end{itemize}

\begin{tcolorbox}[breakable,colback=teal!4,colframe=teal!35!black,title={Before you move on}]
What does \mr{कण्हणे} mean? (\textquotedblleft{}To groan\textquotedblright{}.) Say it once more.
\end{tcolorbox}
\section[\texorpdfstring{\mr{विव्हळणे} (vivhaḷṇe)}{विव्हळणे (vivhaḷṇe)}]{\mr{विव्हळणे} (vivhaḷṇe) — to moan in pain}

\label{lesson:MR-C482-vivhalne}



\begin{quote}
Before the new one: say the Marathi for to get a twinge, then the Marathi for to groan.
\end{quote}

\subsection*{You'll want to know: \mr{विव्हळणे}}
\textbf{\mr{विव्हळणे}} — \emph{vivhaḷṇe} — \textquotedblleft{}to moan in pain\textquotedblright{}.

\begin{quote}
\mr{पाय} \mr{मोडल्यामुळे} \mr{कुत्रा} \mr{विव्हळत} \mr{होता}.
\end{quote}

\begin{quote}
\emph{pāy moḍlyāmuḷe kutrā vivhaḷat hotā.}
\end{quote}

\begin{quote}
The dog was moaning with a broken leg.
\end{quote}

Past (he …): \mr{तो} \mr{विव्हळला}, \emph{to vivhaḷlā}.

Louder than \mr{कण्हणे}.

\begin{cousinweb}
\textbf{\mr{विव्हळणे}} (\emph{vivhaḷṇe}) is usually traced to Sanskrit \emph{vihvala}, \textquotedblleft{}agitated\textquotedblright{} (CDIAL 12038), with the verb ending \emph{-ṇe}; CDIAL gives no Marathi form there.
\end{cousinweb}

\begin{grammarlens}[title={What you've built}]
One more word for describing what you lived and felt.
\end{grammarlens}

\subsection*{Your turn}
\begin{itemize}
\raggedright
  \item \emph{Say it:} \emph{vivhaḷṇe}
  \item \emph{Say it:} \emph{vivhaḷṇe}, once more
  \item \emph{Say it:} \emph{pāy moḍlyāmuḷe kutrā vivhaḷat hotā.} — The dog was moaning with a broken leg.
  \item \emph{Recall:} say the Marathi for to get a twinge, then the Marathi for to groan, then say \emph{vivhaḷṇe} again
\end{itemize}

\begin{tcolorbox}[breakable,colback=teal!4,colframe=teal!35!black,title={Before you move on}]
What does \mr{विव्हळणे} mean? (\textquotedblleft{}To moan in pain\textquotedblright{}.) Say it once more.
\end{tcolorbox}
